\begin{table}[H]
\centering
\sffamily
\small
\renewcommand{\arraystretch}{1.15}
\begin{threeparttable}
\caption{Robustesse internationale : croissance agr\'eg\'ee de la productivit\'e du travail et de la PTF selon la Penn World Table, 2001--2019 (\%)}
\label{tab:note_pwt_robustness}
\begin{tabular}{l *{4}{r}}
\toprule
& $\Delta \ln(Y/L)$ & Rang & PTF agr. & Rang \\
\midrule
Estonie & 4.65 & 1 & 1.40 & 4 \\
Lettonie & 4.10 & 2 & 1.94 & 2 \\
Lituanie & 3.91 & 3 & 2.01 & 1 \\
Tch\'equie & 2.60 & 4 & 1.45 & 3 \\
Danemark & 2.26 & 5 & 0.19 & 10 \\
Autriche & 1.90 & 6 & 0.04 & 11 \\
Espagne & 1.86 & 7 & -0.27 & 18 \\
Slov\'enie & 1.80 & 8 & 0.95 & 5 \\
\'Etats-Unis & 1.53 & 9 & 0.58 & 6 \\
Allemagne & 1.45 & 10 & 0.32 & 8 \\
Su\`ede & 1.43 & 11 & 0.51 & 7 \\
France & 1.33 & 12 & -0.14 & 16 \\
Belgique & 1.18 & 13 & -0.21 & 17 \\
Pays-Bas & 1.11 & 14 & -0.07 & 15 \\
Finlande & 1.09 & 15 & -0.04 & 14 \\
Royaume-Uni & 1.02 & 16 & 0.03 & 12 \\
\textbf{Canada} & \textbf{0.92} & \textbf{17} & \textbf{-0.00} & \textbf{13} \\
Italie & 0.79 & 18 & -0.95 & 19 \\
Japon & 0.40 & 19 & 0.20 & 9 \\
\bottomrule
\end{tabular}
\begin{tablenotes}\footnotesize
\item \textit{Note}\,: M\^eme \'echantillon de 19 pays que le tableau~\ref{tab:note_international}. Croissance annuelle moyenne en points de pourcentage, 2001--2019. $Y/L$ = PIB r\'eel par heure travaill\'ee, calcul\'e comme $rgdpo/(emp \times avh)$. PTF agr. = croissance de $rtfpna$. Les rangs sont calcul\'es sur cet \'echantillon. Cette mesure de PTF est agr\'eg\'ee ; elle n'est donc pas identique \`a la PTF intra-industries du tableau~\ref{tab:note_international}. Source\,: Penn World Table 10.01.
\end{tablenotes}
\end{threeparttable}
\end{table}